% TOP_PROBLEM: {"id": 3419, "title": "Unsolvability of word problem for 2-knot complements", "category_id": 7, "category": {"id": 7, "name": "topology", "display_name": "Topology", "description": "Properties preserved under continuous deformations.", "slug": "topology", "order_index": 7, "created_at": "2026-07-31T15:26:25.671Z"}, "status": "open", "problem_number": "OPG-37237", "set_id": 12}
% FIRST_POSTED: 2026-10-02
% ATTEMPT_STATUS: unresolved
% UnsolvedMath ID 3419; source catalogue code OPG-37237
% !TeX program = lualatex
% GPT-6 Astra Ultra research attempt, 2026-10-02; independently checked partial work.
% Completion estimate: no numerical percentage assigned; see final status and literature audit.
\documentclass[11pt,a4paper]{article}
\usepackage[margin=25mm]{geometry}
\usepackage{fontspec}
\setmainfont{lmroman10-regular.otf}[BoldFont=lmroman10-bold.otf,
 ItalicFont=lmroman10-italic.otf,BoldItalicFont=lmroman10-bolditalic.otf]
\usepackage{amsmath,amssymb,mathtools}
\usepackage{unicode-math}
\setmathfont{latinmodern-math.otf}
\AtBeginDocument{\let\setminus\smallsetminus}
\usepackage{xurl,hyperref}
\hypersetup{colorlinks=true,urlcolor=blue}
\setcounter{secnumdepth}{0}
\setlength{\parindent}{0pt}
\setlength{\parskip}{5pt}
\setlength{\emergencystretch}{3em}
\begin{document}
\section*{Unsolvability of word problem for 2-knot complements}\label{3419}
{\small UnsolvedMath ID 3419 · OPG-37237 · Topology · OpenGarden}
\subsection{Definitions and mathematical statement}
A smooth two-knot is a smooth embedding $e:S^2\hookrightarrow S^4$
in the standard smooth four-sphere. In the PL version, require a
locally flat piecewise-linear embedding: near each point the pair
looks like $(\mathbb R^4,\mathbb R^2)$. Let
$G=\pi_1(S^4\setminus e(S^2))$. Equivalently one may take the
fundamental group of the compact exterior obtained by removing
the interior of a tubular neighborhood. It has a finite presentation
\[
                   G=\langle g_1,\ldots,g_r\mid R_1,\ldots,R_s\rangle.
\]
The word problem for this fixed group asks, on input any finite
word in $g_1^{\pm1},\ldots,g_r^{\pm1}$, whether it represents the
identity. Unsolvability means no algorithm halts and answers correctly
on every such word. This property is independent of the chosen
finite presentation.

Does there exist a smooth two-knot, or a locally flat PL two-knot
in the corresponding version, whose group has unsolvable word
problem? The target is one fixed knot group with this property,
not merely the absence of a uniform algorithm when arbitrary
presentations of different groups are supplied. A construction must
certify the embedding in $S^4$ and establish undecidability for
its complement group. Arbitrary subsets or higher-genus surfaces
do not meet the two-sphere requirement.
\subsection{Short English statement}
Can the complement of a smooth or locally flat PL two-sphere in the four-sphere have a fundamental group with undecidable word problem?
\subsection{Research attempt}
\paragraph{Scope and process.}
GPT-6 Astra Ultra, 2 October 2026. The canonical formulation above is retained. This attempt uses a primary-source literature audit, independent restricted proofs, and adversarial checks. Reproved results are not claimed as new. Numerical tests below are reproducible in \texttt{scripts/check\_attempt\_batch03\_extremal\_2026\_10\_02.py}.

\paragraph{Literature and quantifiers.}
González-Acuña, Gordon and Simon [S3] prove undecidability results about classes of higher-dimensional knot groups and explicitly leave the existence of a 2-knot group with unsolvable word problem as a question. Their results in higher dimension do not realize the required group as a complement of $S^2$ in $S^4$. The target here requires one fixed knot and one fixed finitely presented group whose word problem is undecidable; undecidability of recognizing knots among arbitrary presentations is a different statement. We derive necessary constraints and exclude a broad algorithmically controlled class.

\paragraph{Two necessary group constraints.}
For a smooth or locally flat PL 2-knot $K$, let $E$ be its compact exterior. Its inclusion in $S^4\setminus K$ is a homotopy equivalence, by pushing points in the punctured tubular neighbourhood radially to the boundary. Alexander duality applied to the embedded sphere gives
\[
 H_1(S^4\setminus K;\mathbb Z)\cong H^2(S^2;\mathbb Z)\cong\mathbb Z.
\]
Thus the abelianization of $G=\pi_1(E)$ must be infinite cyclic.

The normal disk bundle of the sphere is trivial: its Euler number is its self-intersection, zero in $S^4$. Hence $S^4$ is obtained by gluing $E$ to $S^2\times D^2$ along $S^2\times S^1$. Van Kampen identifies its fundamental group with the quotient of $G$ by the normal closure of a meridian $\mu$, because the added piece is simply connected and the overlap group is generated by that meridian. Since $S^4$ is simply connected,
\[
 G/\langle\!\langle\mu\rangle\!\rangle=1.
\]
So the group is normally generated by one element. These are necessary conditions; no realization theorem for arbitrary groups satisfying them in dimension four is inferred.

\paragraph{A terminating word algorithm for residual finiteness.}
Let $G=\langle x_1,\ldots,x_d\mid r_1,\ldots,r_s\rangle$ be any fixed finitely presented residually finite group. For an input word $w$, run two effective searches in parallel. First enumerate all finite products of conjugates of $r_j^{\pm1}$ in the free group, freely reducing and comparing with $w$. This search terminates exactly when $w=1$ in $G$, since the defining kernel is the normal closure of the relators.

Second enumerate finite multiplication tables. Retain those satisfying the finite group axioms, enumerate all assignments of the generators to their elements, and retain assignments for which all relators evaluate to the identity. Each retained assignment defines a homomorphism from $G$ to a finite group. If one sends $w$ to a nonidentity element, return nontrivial. Residual finiteness guarantees such a finite quotient whenever $w\ne1$ in $G$. Dovetailing the two searches therefore terminates on every word and gives the correct answer.

This establishes that any desired 2-knot group with undecidable word problem must fail residual finiteness. It is a nonuniform statement for each fixed presentation, which is exactly enough for the necessary condition. It neither assumes nor proves that all 2-knot groups are residually finite, and provides no complexity bound for the searches.

\paragraph{A complete geometric boundary case.}
The unknotted sphere can be viewed inside the decomposition
\[
 S^4=\partial(D^3\times D^2)
 =(S^2\times D^2)\cup(D^3\times S^1),
\]
with corners smoothed. Its exterior is $D^3\times S^1$, and its group is $\mathbb Z$. In the one-generator presentation a word is trivial exactly when its total signed exponent is zero, giving a linear-time bit algorithm by tallying positive and negative occurrences in two binary counters and comparing them at the end. Each counter only increases, so its total carry work is linear in the number of increments. It satisfies both constraints above, showing that those constraints alone do not force undecidability.

\paragraph{Verification and first gap.}
The script checks free reduction and exponent sums on finite sample words, while the dovetailing proof checks both positive and negative termination without relying on experiments. A finitely presented group with an undecidable word problem cannot simply be declared a 2-knot group: an embedding in the correct ambient four-sphere is still required. The missing step is a fixed non-residually-finite undecidable group together with a verified smooth or locally flat PL 2-knot realizing it, or a theorem excluding all such realizations.

\paragraph{Final status.}
The meridian, abelianization and residual-finiteness restrictions are established, and the unknot case is decided. The requested existence statement is \textbf{UNRESOLVED} by this attempt.

\subsection{Sources}
\begin{itemize}
\item[S1] ULAM.AI and the UnsolvedMath Contributors, supplied problems.json, record 3419 (OPG-37237), OpenGarden collection. Catalogue identity and source transcription; CC BY 4.0.
\url{https://huggingface.co/datasets/ulamai/UnsolvedMath}.
\item[S2] Open Problem Garden, Unsolvability of word problem for 2-knot complements. Original problem and discussion; the mathematical formulation below is expanded from the supplied source record.
\url{http://www.openproblemgarden.org/op/unsolvability_of_word_problem_for_2_knot_complements}.
\item[S3] F. González-Acuña, C. McA. Gordon and J. Simon, Unsolvable problems about higher-dimensional knots and related groups, L'Enseignement Mathématique 56 (2010), 143--171; arXiv:0908.4009. See the concluding questions after Theorem 7.1 for the 2-knot distinction.
\url{https://arxiv.org/abs/0908.4009}.
\url{https://ems.press/content/serial-article-files/44199}.
\end{itemize}
\end{document}
